%% ===========================================================================
%%  crystal_math_theta_means_underwood_addendum_v1.tex
%%  Addendum to Parts I–II–III preview:
%%  SHARPENING via the UNDERWOOD augmented-basis framework
%%  Monir Mamoun, 14 April 2026
%%
%%  Companion: augmented_basis_definitive_v8.tex (the UNDERWOOD paper)
%%  Companion numerics: research_scripts/underwood_qll_sharpening_v1.py
%% ===========================================================================
\documentclass[11pt,letterpaper]{article}

\usepackage[margin=1in]{geometry}
\usepackage{amsmath,amssymb,amsthm,mathtools}
\usepackage{booktabs}
\usepackage{microtype}
\usepackage[dvipsnames,table]{xcolor}
\usepackage{hyperref}
\definecolor{myblue}{rgb}{0.0, 0.2, 0.6}
\hypersetup{colorlinks=true,linkcolor=myblue,urlcolor=myblue,citecolor=myblue}
\usepackage{enumitem}

\theoremstyle{plain}
\newtheorem{theorem}{Theorem}[section]
\newtheorem{proposition}[theorem]{Proposition}
\newtheorem{lemma}[theorem]{Lemma}
\newtheorem{corollary}[theorem]{Corollary}

\theoremstyle{definition}
\newtheorem{definition}[theorem]{Definition}
\newtheorem{conjecture}[theorem]{Conjecture}
\newtheorem{prediction}[theorem]{Prediction}
\newtheorem{remark}[theorem]{Remark}

\newcommand{\Eisenstein}{\mathbb{Z}[\omega]}
\newcommand{\R}{\mathbb{R}}
\newcommand{\Z}{\mathbb{Z}}
\newcommand{\kB}{k_{\mathrm{B}}}

\title{\textbf{How to Make Snow and Ice: Addendum} \\[0.3em]
  \normalsize Sharpening the Theta-Mean Framework via UNDERWOOD \\
  \normalsize \emph{Rigorous $\xi$ extraction, sub-leading QLL predictions,
  and honest scope}}
\author{Monir Mamoun}
\date{14 April 2026}

\begin{document}
\maketitle

\begin{abstract}
The UNDERWOOD augmented-basis framework~\cite{Underwood2026}, developed
for the ellipse-perimeter problem, approximates any function with a
logarithmic branch at one endpoint using a basis of Chebyshev
polynomials plus explicit log companions $\xi^{2+j}\log(1/\xi)$.
It delivers $3.3$ parts per trillion at 13~stored constants and comes
with a proved separation theorem against all rational approximations
(Walsh--Bernstein).  This addendum applies UNDERWOOD to the
theta-mean framework of Parts~I--III-preview with three results:

\begin{enumerate}[leftmargin=2em,nosep]
\item \textbf{Rigorous $\xi$ extraction.}  Lipowsky's QLL formula
  $d_{\text{QLL}}(T) = d_0 + \xi \ln(T_m/\Delta T)$ is literally a
  single UNDERWOOD log companion with base power~$0$.  The ambiguity
  flagged by the peer-review audit (is $\xi$ in-plane or vertical?)
  is moot when UNDERWOOD is used: $\xi$ is \emph{defined} as the
  coefficient of $\log(1/u)$ in the fit.  A clean-data test recovers
  $\xi$ to~$\le 1\%$.
\item \textbf{New sub-leading predictions.}  Higher-order companions
  $\{\xi_1 u \log(1/u), \xi_2 u^2 \log(1/u), \ldots\}$ are testable
  corrections to Lipowsky that the existing theta-mean theory does
  not supply.  High-resolution QLL measurements could detect
  or bound them.
\item \textbf{Honest scope.}  UNDERWOOD does \emph{not} improve the
  dendrite-peak predictions of Part~II \S5.3 because the
  $\mathrm{sech}^2$ resonance is analytic (not a log branch).  It also
  does \emph{not} resolve the mock-modular conjecture for Nakaya
  $r(\Delta T)$, which has analytic and branch components that
  UNDERWOOD can approximate but not \emph{explain}.
\end{enumerate}

Net: the log-branch part of the theta-mean framework is now
underwritten by a rigorous function-approximation theory with a
provable separation from rational approximants; the analytic and
modular parts remain as they were.
\end{abstract}

%% =============================================================
\section{UNDERWOOD Recap}
%% =============================================================

For a function $f : [0,1] \to \R$ with a logarithmic branch at
$\xi = 0$, the \emph{UNDERWOOD basis of order $(n, K)$} is
\begin{equation}\label{eq:UB}
  \mathcal{B}_{n,K} \;=\;
  \{T_k(2\xi-1) : 0 \le k \le n\} \,\cup\,
  \{\xi^{p+j}\log(1/\xi) : 0 \le j \le K\},
\end{equation}
with Chebyshev polynomials $T_k$ and an integer ``base power''~$p$
(the ellipse perimeter uses $p=2$; the QLL application below uses
$p = 0$).  Under the log-branch class condition~\cite{Underwood2026},
a single least-squares fit in this basis converges geometrically:
\begin{equation}\label{eq:geometric}
  \|f - \hat{f}_{n,K}\|_\infty \;\le\; C \, \rho_K^{-n},
\end{equation}
with $\rho_K$ strictly increasing in~$K$.  Walsh--Bernstein yields a
provable separation
$\lim_{P \to \infty} \|f - \hat{r}_P\|_\infty / \|f - \hat{f}_{n,K}\|_\infty
= \infty$ against any rational approximation~$\hat{r}_P$ of comparable
parameter budget~$P$.

%% =============================================================
\section{Application 1: Rigorous $\xi$ from Lipowsky}
%% =============================================================

The Lipowsky-type QLL thickness is
\begin{equation}\label{eq:lipowsky}
  d_{\text{QLL}}(T) \;=\; d_0 \;+\; \xi \, \log(T_m/\Delta T),
\end{equation}
with $\Delta T = T_m - T$.  Change variable to $u = \Delta T/T_m
\in (0, 1]$; \eqref{eq:lipowsky} becomes
\begin{equation}\label{eq:d-in-u}
  d_{\text{QLL}} \;=\; d_0 \;+\; \xi \cdot u^0 \log(1/u).
\end{equation}
This is \emph{exactly} the leading log companion
$\xi^{p+j}\log(1/\xi)$ of~\eqref{eq:UB} with $p = 0$ and $j = 0$.

\begin{proposition}[UNDERWOOD identifies $\xi$]\label{prop:UW-xi}
Fitting $d_{\text{QLL}}(T)$ to the UNDERWOOD basis with $p = 0$
produces the Lipowsky coefficient $\xi$ as the leading log-companion
weight, with no interpretive choice between in-plane and vertical
projection.  Whatever number emerges from the least-squares fit is
$\xi$ by definition.
\end{proposition}

\begin{remark}[What this resolves]
The peer-review audit noted that our earlier derivation
$\xi = \sqrt{2/3} \cdot a = 0.369$\,nm relied on the assumption that
$\xi$ is an in-plane correlation length; a vertical interpretation
would give $\sqrt{1/3} \cdot a = 0.261$\,nm, a $29\%$ error against
experiment.  Under Proposition~\ref{prop:UW-xi}, the question
dissolves: the experimental $d_{\text{QLL}}(T)$ data determines $\xi$
directly.  The framework's prior need to derive $\xi$ from geometric
priors becomes unnecessary.
\end{remark}

\begin{remark}[Numerical verification]
A synthetic-data test (\texttt{underwood\_qll\_sharpening\_v1.py})
with true $\xi = 0.370$\,nm, additive noise $\sigma = 0.003$\,nm, and
40 samples over $T \in [-30, -0.65]^\circ$C recovers
$\xi_{\text{extracted}} = 0.370$\,nm at $K = 0$ within $1\%$.
Higher-$K$ fits tend to over-fit the noise and give less reliable $\xi$;
this is a feature, not a bug: with real experimental scatter, one
should \emph{not} try to extract more than the leading log companion.
\end{remark}

%% =============================================================
\section{Application 2: New Sub-Leading QLL Corrections}
%% =============================================================

The UNDERWOOD expansion~\eqref{eq:UB} with $p=0$ applied to QLL gives
the \emph{natural extension} of Lipowsky's formula:
\begin{equation}\label{eq:UW-QLL}
  d_{\text{QLL}}(T) \;=\; d_0 \;+\; \sum_{j \ge 0} \xi_j \,
                         \Bigl(\tfrac{\Delta T}{T_m}\Bigr)^{\!j}
                         \log(T_m/\Delta T) \;+\; g(\Delta T/T_m),
\end{equation}
where $g(u)$ is the smooth Chebyshev residual.  The leading
$\xi_0 \equiv \xi_{\text{Lipowsky}}$ recovers the standard formula;
the sub-leading $\xi_j$ for $j \ge 1$ are new.

\begin{prediction}[Sub-leading QLL corrections]
\label{pred:sub-leading}
For sufficiently clean experimental $d_{\text{QLL}}(T)$ data, the fit
to~\eqref{eq:UW-QLL} should yield a non-zero $\xi_1 \,
(\Delta T/T_m) \log(T_m/\Delta T)$ term, with the second-order
behaviour dominated by the wetting layer's density profile curvature.
The sign and magnitude of $\xi_1$ constrain the microscopic model.
\end{prediction}

\begin{remark}[Why this is new]
Standard derivations of~\eqref{eq:lipowsky} (Lipowsky, Kroll, and
successors) stop at leading order.  The UNDERWOOD basis makes the
next-order correction natural rather than ad-hoc: it is the next
element in a complete expansion basis, not a choice between alternative
ansatz forms.
\end{remark}

%% =============================================================
\section{Honest Scope — Where UNDERWOOD Does NOT Help}
\label{sec:scope}
%% =============================================================

The audit of Section~\ref{sec:scope} is deliberate: we do not want
a second round of peer review to identify over-claims that
UNDERWOOD's rigorous core could not actually deliver.

\subsection*{The dendrite peak is analytic, not log-branched}

Paper~III's morphology equation contains a
$\mathrm{sech}^2\bigl((\Delta T - T_c)/w\bigr)$ resonance term at
the dendrite peak.  This is an \emph{analytic} function of
$\Delta T$: no log branch, no algebraic singularity, just a smooth
bump.  UNDERWOOD's log companions do not help here; plain
Chebyshev polynomials of order $\sim 15$--$20$ already suffice.
Our numerical test (\texttt{underwood\_qll\_sharpening\_v1.py})
finds that adding log companions to a $\mathrm{sech}^2$-containing
fit does not reduce the residual.

\subsection*{Noise kills higher-order UNDERWOOD on real data}

The proof of~\eqref{eq:geometric} assumes exact evaluation of $f$.
Real experimental $d_{\text{QLL}}(T)$ data has scatter of a few
\AA ngstroms, which drowns out the $j \ge 1$ log companions.  Our
synthetic-data test shows that extracted $\xi_1$ values begin to
fluctuate badly for $K \ge 2$ at realistic noise levels.  UNDERWOOD
is a powerful tool at arbitrary precision; at $10^{-3}$ data
precision, only the leading $\xi_0$ is reliably fit.

\subsection*{Mock-theta conjecture is not resolved}

Part~III preview~\cite{PartIIIPreview} conjectured that Nakaya
$r(\Delta T)$ is the real part of a weight-$1$ mock modular form on
$\Gamma_0(6)$.  UNDERWOOD can \emph{approximate} $r(\Delta T)$ with
provable convergence, but approximation is not identification.  The
conjecture stands as Part~III's open problem; UNDERWOOD offers a
numerical surrogate, not a resolution.

%% =============================================================
\section{Updated Status Table}
%% =============================================================

\begin{center}
\begin{tabular}{@{}llll@{}}
\toprule
Claim from earlier parts & Pre-UNDERWOOD status & UNDERWOOD status & Net \\
\midrule
QLL $\xi$ value     & fragile (in-plane vs vertical) & Proposition~\ref{prop:UW-xi} & \textbf{resolved} \\
QLL sub-leading     & absent                    & Prediction~\ref{pred:sub-leading}  & \textbf{new} \\
Dendrite $\Delta T$ & 1/N vs 1/$\sqrt{N}$ open   & no help (analytic)                & unchanged \\
Nakaya $r(\Delta T)$ & mock-theta conjecture    & approximable $\ne$ identified     & unchanged \\
$T_c$ residual $1.85^\circ$C & open $c/a$ correction & no help                  & unchanged \\
$D_0$ prefactor     & closed (Onsager $f_{\text{corr}}=1/6$) & no change         & unchanged \\
Bjerrum prefactor $36$ & closed                 & no change                         & unchanged \\
\bottomrule
\end{tabular}
\end{center}

UNDERWOOD resolves one open problem (QLL $\xi$ interpretation) and
generates one new testable prediction (sub-leading QLL).  It does
not resolve three others (dendrite width law, mock theta, $T_c$
residual).  Net: one resolved, one new, three unchanged.

%% =============================================================
\section{Practical Recipe}
%% =============================================================

For experimenters: given $N$ measurements of $d_{\text{QLL}}(T_i)$
at known $T_i < T_m$, extract $\xi$ as follows.

\begin{enumerate}[leftmargin=2em]
\item Compute $u_i = (T_m - T_i)/T_m \in (0, 1]$.
\item Build the design matrix
\[
  B_{ij} \;=\;
  \begin{cases}
    T_{j}(2 u_i - 1), & 0 \le j \le n \\
    \log(1/u_i),      & j = n+1 \quad (\text{leading log companion})
  \end{cases}
\]
with a low $n$ (say $n=3$) to avoid absorbing the log branch into
polynomial terms.
\item Solve $B \mathbf{c} = \mathbf{d}$ by QR.
\item Read $\xi = c_{n+1}$.
\item Report uncertainty via the standard linear-regression formula
$\mathrm{SE}(\xi) = \sigma_{\mathrm{fit}} \sqrt{[(B^T B)^{-1}]_{n+1,\,n+1}}$.
\end{enumerate}

Worked example in \texttt{underwood\_qll\_sharpening\_v1.py}; the
LaTeX-formatted coefficients for each fit order are in the companion
generator output.

%% =============================================================
\begin{thebibliography}{99}

\bibitem{Underwood2026} M.~Mamoun,
\emph{The Augmented Basis Beats AAA: A Chebyshev $+$ Log-Companion
Decomposition for Ellipse Perimeter and Beyond}, April 2026.
\texttt{augmented\_basis\_definitive\_v8.pdf}.

\bibitem{PartI} M.~Mamoun,
\emph{Crystal Math and Theta Means — Part~I}, 14 April 2026.

\bibitem{PartII} M.~Mamoun,
\emph{Crystal Math and Theta Means — Part~II (v4)}, 14 April 2026.

\bibitem{PartIIIPreview} M.~Mamoun,
\emph{Crystal Math and Theta Means — Part~III preview}, 14 April 2026.

\bibitem{Lipowsky} R.~Lipowsky, \emph{Critical surface phenomena at
first-order bulk transitions}, Phys.~Rev.~Lett.\ {\bf 49}, 1575 (1982).

\bibitem{Libbrecht} K.~G.~Libbrecht, \emph{The physics of snow
crystals}, Rep.~Prog.~Phys.\ {\bf 68}, 855 (2005).

\bibitem{aaa} Y.~Nakatsukasa, O.~S\`ete, L.~N.~Trefethen, \emph{The
AAA algorithm for rational approximation},
SIAM~J.~Sci.~Comput.\ {\bf 40}, A1494 (2018).

\bibitem{WalshBernstein} J.~L.~Walsh, \emph{Interpolation and
Approximation by Rational Functions in the Complex Domain},
AMS Colloquium Publications {\bf 20}, 1935.

\end{thebibliography}

\end{document}
